% draft_conclusion.tex
% Style: Academic Paper
% Target: ~500--700 words of body text
% Rewritten 2026-05-17 — shortened and aligned with draft_discussion.tex

\section{Conclusion}\label{conclusion}

This paper examines whether Twitter sentiment causally affects intraday stock
returns, using a Bloomberg panel of S\&P~500 and Russell~3000 firms and two
identification strategies: firm fixed-effects OLS and Double Machine Learning
with XGBoost nuisance learning. The central finding is conditional. Twitter
sentiment does not reliably predict returns in large-cap S\&P~500 firms, where
the null result holds across every fixed-effects specification and every
calendar year from 2016 through 2025. In the broader Russell~3000, sentiment
effects are positive, significant, and increasing in magnitude as market
capitalization decreases---small-cap coefficient $+0.942^{***}$, mid-cap
$+0.804^{***}$, large-cap $+0.375^{***}$. This gradient is consistent with
an information-environment explanation where Twitter commentary carries marginal
information precisely where analyst and institutional coverage is thinnest.

Two additional results support the presence of a real, if modest, signal. The
DoubleML estimator recovers a significant negative coefficient on Twitter
sentiment ($-0.924^{***}$) in the S\&P~500 panel, precisely estimated via
five-fold cross-fitting and confirmed as stable across the robustness specifications. And a long-short portfolio sorted on
daily sentiment earns positive annualized returns in nine of ten complete
calendar years, even in years when the cross-sectional regression coefficient
is near zero---consistent with a nonlinear tail that linear regression
cannot recover. Taken together, these results give cautious confidence that
the sentiment signal exists but require that its practical use take into
account the data and identification limitations described below.

The primary caveat is reverse causality. The placebo test, regressing
yesterday's return on today's sentiment, produces a coefficient of
$+2.348^{***}$, larger in magnitude than the forward DoubleML estimate. This
does not eliminate the possibility of a forward effect, but it establishes that
the sentiment-return correlation is contaminated by prices feeding
back into social media activity. The Bloomberg daily aggregate compounds this
problem by conflating pre-market tweets with same-day price reactions, making
it impossible to cleanly separate cause from effect without timestamp-level
data. Additionally, the inclusion of same-day intraday price controls
(\texttt{px\_high}, \texttt{px\_low}) in the FE-OLS specification introduces a
bad-control problem that likely attenuates the Twitter coefficient toward zero;
DoubleML's nonparametric nuisance learner handles this implicitly, which is the
preferred approach for future work in this area.

This paper makes three contributions. First, the Russell~3000 cap-group
decomposition refines Gu and Kurov (2020): their positive result likely
reflects smaller firms in their universe rather than a universal large-cap
effect, and the monotone cap-group gradient provides a direct test of the
information-environment hypothesis. Second, this is the first application of
the Chernozhukov et al.\ (2018) double machine learning framework to the
Twitter sentiment and stock return problem, demonstrating its viability as a
debiased estimator in high-dimensional sentiment panels. Third, the use of
univariate DoubleML estimates at lags one through seven as a reverse causality
screen---a lag-CATE design---is a methodological contribution transferable to
any setting where a researcher needs to assess whether sentiment leads or lags
prices under nonlinear confounding.

There is meaningful evidence that Twitter sentiment predicts returns for
smaller stocks. To assert causality in that segment, however, these findings
would need to pass the same robustness checks applied here to the S\&P~500: a
cap-group-specific placebo test, lag-decay analysis by size partition, and
within-group DoubleML estimation. Extending those tests, and replacing the
Bloomberg daily aggregate with a time-stamped, entity-level sentiment
instrument, is the subject of future work.
